% =====================================================================
%  Lecture 17 -- 3GPP Standardisation                      (Persian edition)
% =====================================================================
\documentclass[11pt,a4paper]{article}
\providecommand\ACNroot{../..}
\input{\ACNroot/templates/lecture-fa}

\ACNset{number}{17}
\ACNset{title}{استانداردسازی \lr{3GPP}}
\ACNset{shorttitle}{استانداردسازی \lr{3GPP}}
\ACNset{titleother}{3GPP Standardisation}
\ACNset{part}{بخش سوم -- معماری‌های نرم‌افزارمحور (درس‌های ۱۵ تا ۱۸)}
\ACNset{time}{۵ تا ۷ ساعت (۳ ساعت مطالعه، ۲ تا ۴ ساعت تمرین)}
\ACNset{level}{مقدماتی تا متوسط}
\ACNset{prereq}{درس‌های ۱ تا ۱۶ (به‌عنوان نمونه)؛ احتمال مقدماتی}
\ACNset{gen}{5}

\begin{document}
\ACNcover

% ---------------------------------------------------------------------
\section{اهداف یادگیری}
پس از مطالعهٔ این درس باید بتوانید:
\begin{objectives}
  \item سازمان \acro{3GPP} -- شرکای سازمانی، گروه‌های مشخصات فنی و گروه‌های کاری -- و رابطهٔ آن با \lr{ITU-R} و \acro{IETF} را
        \textbf{توصیف کنید}.
  \item روش سه‌مرحله‌ای، شماره‌گذاری مشخصات و مدل انتشار را \textbf{توضیح دهید}.
  \item شمارهٔ نسخهٔ یک مشخصه را \textbf{رمزگشایی کنید} و درخواست‌های تغییر را \textbf{دسته‌بندی کنید}.
  \item مسیر یک قابلیت از مورد مطالعه تا مشخصه و تا محصول را \textbf{دنبال کنید}.
  \item یک مشخصه را به‌طور کارآمد \textbf{بیابید و بخوانید}: دامنه، متن الزام‌آور، پیوست‌ها.
  \item دربارهٔ هم‌کنش‌پذیری، انباشت کار و زمان رسیدن به بازار \textbf{استدلال کمّی کنید}.
\end{objectives}

% ---------------------------------------------------------------------
\section{دانش پیش‌نیاز}
این درس دربارهٔ فرایندی است که پشت همهٔ مطالب این درس‌نامه قرار دارد. درس‌های پیشین نمونه‌ها را فراهم می‌کنند: مشخصاتی که
در آن‌ها ارجاع داده شد (\lr{TS 23.401}، \lr{TS 29.272}، \lr{TS 38.300}، \ldots) محصول فرایندی‌اند که اینجا توصیف می‌شود.
تمرین ۵ از قانون لیتل (درس ۱۵) استفاده می‌کند.

\paragraph{واژه‌های کلیدی} واژگان این درس همراه با معادل انگلیسی آن‌ها:
\begin{keyterms}
پروژهٔ مشارکت نسل سوم & \lr{3GPP} & مشارکت سازمان‌های استاندارد که سامانه‌های موبایل را مشخص می‌کند. \\
شریک سازمانی & \lr{Organisational partner} & نهاد استاندارد منطقه‌ای/ملی که عضو \acro{3GPP} است. \\
گروه مشخصات فنی / گروه کاری & \lr{TSG / WG} & گروه‌های تصمیم‌گیر و زیرگروه‌های فنی آن‌ها. \\
مشخصات فنی / گزارش فنی & \lr{TS / TR} & مشخصهٔ الزام‌آور / گزارش مطالعهٔ آگاهی‌بخش. \\
مراحل ۱، ۲، ۳ & \lr{Stage 1, 2, 3} & نیازمندی‌ها، معماری، پروتکل‌های تفصیلی. \\
انتشار & \lr{Release} & مجموعهٔ سازگاری از قابلیت‌ها و مشخصات که با هم منجمد می‌شوند. \\
انجماد & \lr{Freeze} & نقطه‌ای که پس از آن فقط اصلاحات پذیرفته می‌شوند. \\
درخواست تغییر & \lr{Change request (CR)} & پیشنهاد رسمی برای تغییر مشخصه‌ای تحت کنترل تغییر. \\
مورد کاری / مطالعه & \lr{Work item / study item} & وظیفهٔ مصوب برای مشخص‌کردن / مطالعهٔ یک قابلیت. \\
پتنت ضروری استاندارد & \lr{SEP / FRAND} & پتنت ضروری استاندارد / مجوزدهی منصفانه، معقول و غیرتبعیض‌آمیز. \\
\end{keyterms}

% ---------------------------------------------------------------------
\section{مقدمه و انگیزه}
گوشی‌ای که در یک کشور خریده می‌شود روی شبکه‌ای در کشوری دیگر که از تجهیزات چند فروشنده ساخته شده کار می‌کند. اینکه این
اصلاً کار می‌کند نتیجهٔ استانداردهاست: هزاران سند که هر پیام، هر زمان‌سنج و هر بیت را مشخص می‌کنند و رقبایی آن‌ها را
می‌نویسند که در همان اتاق‌ها همکاری می‌کنند. برای شبکه‌های موبایل نهاد غالب \emph{پروژهٔ مشارکت نسل سوم} (\acro{3GPP}) است
که در ۱۹۹۸ برای مشخص‌کردن نسل سوم ایجاد شد و از آن پس مسئول تکامل \acro{GSM}، \acro{LTE}، نسل پنجم و اکنون نسل ششم بوده
است.

مهندسان صنعت موبایل بخش زیادی از وقت خود را با مشخصات \acro{3GPP} می‌گذرانند: پیاده‌سازی، آزمون در برابر آن‌ها یا مشارکت در
نگارششان. بنابراین فهم سازمان اسناد، نحوهٔ تکامل آن‌ها و نحوهٔ تصمیم‌گیری مهارتی عملی است، نه دانستنی‌ای حاشیه‌ای. این فهم
همچنین توضیح می‌دهد چرا سامانه‌ها چنین‌اند: بسیاری از انتخاب‌های طراحی در این درس‌نامه مصالحه‌هایی‌اند که در جلسات
استانداردسازی به دست آمده‌اند.

% ---------------------------------------------------------------------
\section{مفاهیم بنیادی}

\subsection{سازمان}
\acro{3GPP} خود یک نهاد حقوقی نیست: مشارکتی از هفت \emph{شریک سازمانی} است -- \acro{ARIB} و \acro{TTC} (ژاپن)، \acro{ATIS}
(آمریکا)، \acro{CCSA} (چین)، \acro{ETSI} (اروپا)، \acro{TSDSI} (هند) و \acro{TTA} (کره) -- که مشخصات آن را به استانداردهای خود
منتقل می‌کنند. شرکت‌ها به‌عنوان عضو یکی از این شرکا مشارکت می‌کنند؛ \emph{شرکای نمایندهٔ بازار} (انجمن‌های صنعتی)
نیازمندی‌ها را عرضه می‌کنند. کار فنی در سه \emph{گروه مشخصات فنی} و گروه‌های کاری آن‌ها انجام می‌شود (\cref{fig:org})
\cite{tr21900}:
\begin{itemize}
  \item \textbf{\lr{TSG RAN}}: دسترسی رادیویی (لایهٔ فیزیکی، پروتکل‌های رادیویی، معماری \acro{RAN}، \acro{RF}، آزمون‌های
        انطباق \acro{UE})؛
  \item \textbf{\lr{TSG SA}}: جنبه‌های خدمت و سامانه (نیازمندی‌ها، معماری، امنیت، کدک‌ها، مدیریت) و هماهنگی کلی؛
  \item \textbf{\lr{TSG CT}}: شبکهٔ هسته و پایانه‌ها (\acro{NAS}، پروتکل‌های هسته، سیاست، \acro{USIM}).
\end{itemize}
گروه‌های کاری اسناد را آماده می‌کنند؛ جلسات عمومی \acro{TSG} که هر فصل برگزار می‌شوند آن‌ها را تصویب می‌کنند. تصمیم‌ها با
اجماع گرفته می‌شوند؛ رأی‌گیری رسمی ممکن اما نادر است.

\begin{acnfigure}
  \resizebox{\linewidth}{!}{\acnfig{lecture-17/organisation}}
  \caption{سازمان \acro{3GPP}: شرکای سازمانی، هماهنگی پروژه، \acro{TSG}ها و گروه‌های کاری آن‌ها (گزیده).}
  \label{fig:org}
\end{acnfigure}

\subsection{رابطه با نهادهای دیگر}
\lr{ITU-R} نیازمندی‌های هر نسل \acro{IMT} را تعریف می‌کند (درس ۱۴) \cite{itu2017m2410,itu2023m2160}؛ فناوری‌های \acro{3GPP} از
طریق شرکای سازمانی به‌عنوان واسط‌های رادیویی نامزد \acro{IMT} ارائه و ارزیابی می‌شوند. \acro{IETF} پروتکل‌های اینترنتی را
مشخص می‌کند که \acro{3GPP} دوباره به کار می‌برد -- \acro{SCTP}، \lr{Diameter}، \lr{HTTP/2}، \acro{SIP}، \acro{IPsec}،
\acro{TLS} \cite{rfc6733,rfc9113,rfc9260,rfc3261} -- و \acro{3GPP} در صورت نیاز افزونه درخواست می‌کند. نهادهای صنعتی مانند
\acro{GSMA} رهنمودهای عملیاتی (مثلاً رومینگ، نمایه‌های \acro{VoLTE}) را تعریف می‌کنند و سازمان‌های گواهی‌دهی دستگاه‌ها را
در برابر مشخصات انطباق \acro{3GPP} می‌آزمایند.

\subsection{سه مرحله}
هر قابلیت در سه سطح مشخص می‌شود (\cref{fig:stages}):
\begin{itemize}
  \item \textbf{مرحلهٔ ۱} خدمت را از دید کاربر توصیف می‌کند: نیازمندی‌ها، بدون راه‌حل (مثلاً \lr{TS 22.261} \cite{ts22261}).
  \item \textbf{مرحلهٔ ۲} معماری را توصیف می‌کند: موجودیت‌های کارکردی، نقاط مرجع و جریان‌های اطلاعات (مثلاً \lr{TS 23.501}
        \cite{ts23501}).
  \item \textbf{مرحلهٔ ۳} پروتکل‌ها را بیت به بیت مشخص می‌کند: پیام‌ها، عناصر اطلاعاتی، کدگذاری‌ها، ماشین‌های حالت (مثلاً
        \lr{TS 24.501} \cite{ts24501}، \lr{TS 38.331} \cite{ts38331}).
\end{itemize}
سپس مشخصات آزمون انطباق تعیین می‌کنند دستگاه‌ها چگونه در برابر مرحلهٔ ۳ آزموده شوند.

\begin{acnfigure}
  \resizebox{\linewidth}{!}{\acnfig{lecture-17/stages}}
  \caption{روش سه‌مرحله‌ای با مشخصات نمونه.}
  \label{fig:stages}
\end{acnfigure}

\subsection{شماره‌گذاری مشخصات}
هر مشخصه یک شمارهٔ \lr{ss.nnn} دارد که سری \lr{ss} آن موضوع را نشان می‌دهد \cite{tr21900}.
\begin{table}[htbp]
\centering\small
\caption{برخی سری‌های مشخصات \acro{3GPP}.}\label{tab:series}
\rowcolors{2}{acnmist}{white}
\begin{tabularx}{\linewidth}{@{}l Y l@{}}
\toprule
\rowcolor{white}سری & موضوع & نمونه \\
\midrule
\lr{21} & نیازمندی‌ها، واژگان، روش‌های کار          & \lr{TR 21.905} \cite{tr21905} \\
\lr{22} & جنبه‌های خدمت (مرحلهٔ ۱)                   & \lr{TS 22.261} \\
\lr{23} & تحقق فنی (مرحلهٔ ۲)                        & \lr{TS 23.401, 23.501} \\
\lr{24} & پروتکل‌های سیگنالینگ \lr{UE}--هسته (مرحلهٔ ۳) & \lr{TS 24.301, 24.501} \\
\lr{26} & کدک‌ها                                     & -- \\
\lr{28, 32} & مدیریت و صورت‌حساب                     & \lr{TS 28.530} \\
\lr{29} & پروتکل‌های شبکهٔ هسته (مرحلهٔ ۳)            & \lr{TS 29.274, 29.500} \\
\lr{31} & \lr{SIM/USIM}                              & -- \\
\lr{33} & امنیت                                      & \lr{TS 33.401, 33.501} \\
\lr{35} & الگوریتم‌های امنیتی                        & \lr{TS 35.206} \\
\lr{36} & رادیوی \acro{LTE} (\lr{E-UTRA})             & \lr{TS 36.300, 36.331} \\
\lr{37} & چند فناوری رادیویی                         & \lr{TS 37.340} \\
\lr{38} & رادیوی \acro{NR}                            & \lr{TS 38.300, 38.331} \\
\bottomrule
\end{tabularx}
\end{table}
یک \emph{مشخصهٔ فنی} (\acro{TS}) الزام‌آور است: محصولات باید با آن مطابقت داشته باشند. یک \emph{گزارش فنی} (\acro{TR})
آگاهی‌بخش است: مطالعات، ارزیابی‌ها و گزینه‌ها را ثبت می‌کند، اغلب در بازه‌های \lr{8xx} و \lr{9xx} (مثلاً \lr{TR 38.913}
\cite{tr38913}، \lr{TR 38.901} \cite{tr38901}).

% ---------------------------------------------------------------------
\section{توضیحات فنی تفصیلی}

\subsection{انتشارها}
\acro{3GPP} کار خود را در \emph{انتشارها} منتشر می‌کند: مجموعه‌های سازگاری از قابلیت‌ها که با هم مشخص می‌شوند
(\cref{fig:releases}). هر انتشار تاریخ‌های انجماد دارد: پس از \emph{انجماد مرحلهٔ ۳} هیچ کارکرد جدیدی افزوده نمی‌شود و فقط
اصلاحات پذیرفته می‌شوند؛ پس از \emph{انجماد \lr{ASN.1}} کدگذاری پیام‌های پروتکل رادیویی پایدار است تا پیاده‌سازی‌ها
هم‌کنش‌پذیر شوند. انتشارها هم‌پوشانی دارند: در حالی که یکی اصلاح می‌شود، بعدی در حال توسعه و بعد از آن در حال مطالعه است.
فروشندگان و اپراتورها زیرمجموعه‌ای از قابلیت‌های یک انتشار را پیاده می‌کنند و دستگاه‌ها قابلیت‌های خود را اعلام می‌کنند،
پس «پشتیبانی از انتشار ۱۷» بدون فهرست قابلیت‌ها معنای چندانی ندارد.

\begin{acnfigure}
  \resizebox{\linewidth}{!}{\acnfig{lecture-17/releases}}
  \caption{برخی انتشارهای \acro{3GPP} با قابلیت‌های شاخص. سال‌ها زمان تقریبی تکمیل را نشان می‌دهند؛ رنگ فیروزه‌ای انتشارهای نسل پنجم است.}
  \label{fig:releases}
\end{acnfigure}

\subsection{نسخه‌ها}
هر مشخصه یک نسخهٔ \lr{x.y.z} دارد \cite{tr21900}. پیش از تصویب، $x = 1$ (ارائه برای اطلاع) یا $x = 2$ (ارائه برای تصویب)
است. پس از تصویب و قرارگرفتن تحت کنترل تغییر، $x$ برابر شمارهٔ انتشار است (مثلاً ۱۷ برای انتشار ۱۷)؛ $y$ با هر مجموعهٔ
تغییرات فنی مصوب افزایش می‌یابد؛ $z$ برای تغییرات صرفاً ویرایشی افزایش می‌یابد. بنابراین یک مشخصه در نسخه‌های موازی، یکی
برای هر انتشار، وجود دارد و اصلاحات اغلب در چند انتشار بازتاب داده می‌شوند.

\subsection{درخواست‌های تغییر}
مشخصه‌ای که تحت کنترل تغییر است فقط از طریق \emph{درخواست‌های تغییر} (\acro{CR}) تغییر می‌کند که هر یک دسته‌ای دارد
\cite{tr21900}: \textbf{\lr{F}} (اصلاح یک خطا)، \textbf{\lr{A}} (بازتاب یک اصلاح دستهٔ \lr{F} در انتشار بعدی)،
\textbf{\lr{B}} (افزودن یک قابلیت)، \textbf{\lr{C}} (تغییر کارکردی یک قابلیت) و \textbf{\lr{D}} (ویرایشی). پس از انجماد معمولاً
فقط \lr{F}، \lr{A} و \lr{D} پذیرفته می‌شوند. \acro{CR}ها در گروه کاری توافق و در \acro{TSG} تصویب می‌شوند؛ سپس مشخصه با نسخهٔ
افزایش‌یافته دوباره منتشر می‌شود (\cref{fig:cr}).

\begin{acnfigure}
  \resizebox{\linewidth}{!}{\acnfig{lecture-17/cr-process}}
  \caption{از مورد مطالعه تا مشخصه، و چرخهٔ درخواست تغییر.}
  \label{fig:cr}
\end{acnfigure}

\subsection{نحوهٔ انجام کار}
شرکت‌ها مشارکت‌های مکتوب (اسناد موقت، «\lr{Tdoc}») را به جلسات گروه کاری ارائه می‌دهند که چند بار در هر چرخهٔ انتشار در
مناطق مختلف و با جلسات الکترونیکی در میانه برگزار می‌شوند. پیشنهادها بحث، ادغام و بازنگری می‌شوند تا اجماع پدید آید؛
رئیسان توافق‌ها را ثبت می‌کنند. مشارکت‌های معمول عبارت‌اند از پیشنهاد راه‌حل در یک مطالعه، پیشنهاد متن برای یک \acro{TR} یا
\acro{TS} و \acro{CR}ها. همهٔ اسناد در وب‌گاه \acro{3GPP} عمومی‌اند و این تاریخچهٔ هر تصمیم طراحی را ردیابی‌پذیر می‌کند.

\subsection{خواندن یک مشخصه}
مشخصهٔ \acro{3GPP} از الگویی ثابت پیروی می‌کند: دامنه، مراجع، تعاریف و اختصارات، سپس بندهای فنی و پیوست‌ها. متن الزام‌آور از
«\lr{shall}» (اجباری)، «\lr{should}» (توصیه‌شده) و «\lr{may}» (مجاز) استفاده می‌کند؛ پیوست‌های آگاهی‌بخش پیش‌زمینه و نمونه
می‌دهند. توصیهٔ عملی: نخست توصیف مرحلهٔ ۲، سپس رویهٔ مرحلهٔ ۳ و تنها پس از آن تعریف پیام‌ها و عناصر اطلاعاتی را بخوانید؛ انتشار
و نسخه را بررسی کنید؛ و در میان مشخصات جست‌وجو کنید، چون یک رویه معمولاً میان چند مشخصه تقسیم شده است (مثلاً یک اتصال
\lr{TS 23.401}، \lr{24.301}، \lr{36.413}، \lr{29.274} و \lr{29.272} را درگیر می‌کند).

\subsection{مالکیت فکری}
شرکت‌هایی که فناوری عرضه می‌کنند اغلب پتنت‌هایی دارند که برای پیاده‌سازی یک استاندارد ضروری‌اند (\emph{پتنت‌های ضروری
استاندارد}، \acro{SEP}). بر اساس سیاست‌های مالکیت فکری شرکای سازمانی، آن‌ها این پتنت‌ها را اعلام و تعهد می‌کنند با شرایط
منصفانه، معقول و غیرتبعیض‌آمیز (\acro{FRAND}) مجوز دهند. مجوزدهی به‌صورت دوجانبه و بیرون از \acro{3GPP} مذاکره می‌شود و
سطح کل حق امتیازها بحثی همیشگی است (\cref{sec:stack}).

% ---------------------------------------------------------------------
\section{مدل‌های ریاضی}
مدل‌های زیر عمداً ساده‌اند؛ مثال‌های عددی آن‌ها از فرض‌های تصویری استفاده می‌کنند، نه آمار اندازه‌گیری‌شدهٔ \acro{3GPP}.

\subsection{هم‌کنش‌پذیری در میان واسط‌ها}\label{sec:interop}
\paragraph{مدل} یک خدمت سرتاسری از $k$ واسط چندفروشنده‌ای می‌گذرد. اگر هر یک با احتمال $p$ و مستقل از دیگران درست پیاده
شده باشد، خدمت با این احتمال کار می‌کند:
\begin{equation}\label{eq:interop}
  P_{\mathrm{e2e}} = p^{k}.
\end{equation}
\paragraph{مثال} $p = 0.98$ و $k = 20$: $P_{\mathrm{e2e}} = 0.67$. رسیدن به ۰٫۹ با $k = 20$ مستلزم $p \ge 0.9^{1/20} = 0.995$
است.
\paragraph{تفسیر} هرچه سامانه واسط‌های بیشتری داشته باشد، مشخصات باید دقیق‌تر و آزمون کامل‌تر باشد -- دلیل جزئیات مرحلهٔ ۳،
آزمون‌های انطباق و رویدادهای آزمون هم‌کنش‌پذیری. این همچنین استدلالی به سود واسط‌های کمتر و ساده‌تر است.

\subsection{کار در جریان}\label{sec:little}
\paragraph{مدل} اگر \acro{CR}های یک گروه مشخصات با نرخ $\lambda$ (در ماه) برسند و هر یک به‌طور میانگین $T$ ماه از نخستین ارائه
تا تصویب بماند، قانون لیتل تعداد \acro{CR}های باز را می‌دهد:
\begin{equation}\label{eq:little}
  N = \lambda\,T .
\end{equation}
\paragraph{مثال} با $\lambda = 400$ \acro{CR} در ماه و $T = 3$ ماه، در هر لحظه حدود ۱۲۰۰ \acro{CR} باز است. نصف‌کردن $T$ (مثلاً
با جلسات الکترونیکی میان جلسات عمومی) انباشت را با همان گذردهی نصف می‌کند.

\subsection{زمان رسیدن به بازار}\label{sec:ttm}
\paragraph{مدل} یک قابلیت از مطالعه، کار الزام‌آور، پیاده‌سازی و آزمون می‌گذرد؛ مراحل می‌توانند تا حدی هم‌پوشانی داشته
باشند. زمان از آغاز مطالعه تا در دسترس بودن تجاری
\begin{equation}\label{eq:ttm}
  T_{\mathrm{market}} = \sum_{i} T_i - \sum_{j} \Delta_j
\end{equation}
است، با طول مراحل $T_i$ و هم‌پوشانی‌های $\Delta_j$.
\paragraph{مثال} مطالعه ۱۲ ماه، مرحلهٔ ۲ دوازده، مرحلهٔ ۳ دوازده، پیاده‌سازی ۱۸ و آزمون و گواهی ۶ ماه، با ۶ ماه هم‌پوشانی میان
پیاده‌سازی و انتهای مرحلهٔ ۳: $T_{\mathrm{market}} = 54$ ماه، یعنی حدود ۴٫۵ سال -- سازگار با فاصلهٔ چندساله میان آغاز مطالعات
یک نسل و نخستین شبکه‌های تجاری آن.

\subsection{انباشت حق امتیاز}\label{sec:stack}
\paragraph{مدل} قیمت یک دستگاه $P$ حامل حق امتیاز برای $N$ دارندهٔ \acro{SEP} است. اگر هر یک مستقلاً نرخ $r_i$ بخواهد، مجموع
$R = \sum_i r_i$ است؛ با سقف کل $R_{\max}$ که بر اساس سهم $s_i$ هر دارنده از پتنت‌های ضروری تقسیم شود،
\begin{equation}\label{eq:stack}
  R = \sum_{i=1}^{N} r_i, \qquad r_i^{\mathrm{cap}} = R_{\max}\, s_i, \quad \sum_i s_i = 1 .
\end{equation}
\paragraph{مثال} بیست دارنده که هر کدام ۱٪ بخواهند به ۲۰٪ قیمت دستگاه می‌رسند؛ با سقف ۵٪، دارنده‌ای با ۱۰٪ پتنت‌های ضروری
۰٫۵٪ دریافت می‌کند. اقتصاد \acro{SEP}ها محل مناقشه است؛ مدل فقط نشان می‌دهد چرا این پرسش مطرح می‌شود.

% ---------------------------------------------------------------------
\section{معماری شبکه و نمودارها}
سازمان \acro{3GPP} بازتاب معماری سامانه‌هایی است که مشخص می‌کند: گروه‌های کاری \acro{RAN} مالک واسط رادیویی و واسط‌های
\acro{RAN} (\lr{Uu}، \lr{X2/Xn}، \lr{S1/NG}، \lr{F1}) هستند، \lr{SA2} مالک معماری سامانه و نقاط مرجع، و گروه‌های \acro{CT}
مالک پروتکل‌های هسته و \acro{NAS} میان \acro{UE} و هسته. قابلیت جدیدی که از این مرزها می‌گذرد -- مثلاً برش‌بندی شبکه -- به
موارد کاری هماهنگ در چند گروه نیاز دارد که مرحلهٔ ۲ \lr{SA2} مرجع مشترک آن‌هاست.

% ---------------------------------------------------------------------
\section{رویه‌های گام‌به‌گام}

\begin{procedurebox}{معرفی یک قابلیت جدید}\label{proc:feature}
\begin{steps}
  \item شرکت‌ها یک مورد مطالعه (\acro{SI}) با اهداف پیشنهاد می‌کنند؛ \acro{TSG} آن را تصویب می‌کند.
  \item گروه‌های کاری موارد کاربرد و راه‌حل‌های نامزد را مطالعه و در یک \acro{TR} با ارزیابی‌ها و نتیجه‌گیری مستند می‌کنند.
  \item یک مورد کاری (\acro{WI}) با دامنه، انتشار هدف و گروه‌های مسئول تصویب می‌شود.
  \item مشخصات مرحلهٔ ۲ و ۳ نوشته یا گسترش داده می‌شوند (\acro{TS} جدید یا \acro{CR} بر مشخصات موجود) تا انجمادهای مرحلهٔ ۳ و
        \lr{ASN.1}.
  \item پس از انجماد، اصلاحات از طریق \acro{CR}های دستهٔ \lr{F} ادامه می‌یابد؛ آزمون‌های انطباق مشخص می‌شوند؛ محصولات قابلیت
        را پیاده و گواهی می‌کنند.
\end{steps}
\end{procedurebox}

\begin{procedurebox}{یافتن مشخصهٔ درست}\label{proc:find}
\begin{steps}
  \item موضوع و مرحله را مشخص کن: نیازمندی (۲۲)، معماری (۲۳)، پروتکل (۲۴، ۲۹، ۳۶، ۳۸)، امنیت (۳۳).
  \item از نمای کلی مرحلهٔ ۲ (مثلاً \lr{TS 23.501} یا \lr{TS 38.300}) آغاز کن و مراجع آن را دنبال کن.
  \item انتشاری را که با پیاده‌سازی تطابق دارد برگزین و آخرین نسخهٔ آن انتشار را بارگیری کن.
  \item دامنه و تعاریف، سپس بند رویه، سپس تعریف پیام‌ها را بخوان؛ تاریخچهٔ \acro{CR} را برای تغییرات اخیر بررسی کن.
\end{steps}
\end{procedurebox}

% ---------------------------------------------------------------------
\section{مثال‌های حل‌شده}

\begin{examplebox}{رمزگشایی یک نسخه}\label{ex:version}
\textbf{مسئله.} «\lr{TS 38.331 V17.5.0}» و «\lr{TS 23.501 V2.0.0}» چه می‌گویند؟

\textbf{حل.} \lr{TS 38.331} در سری ۳۸ (رادیوی \acro{NR}) است: پروتکل \acro{RRC} در \acro{NR}. \lr{V17.5.0}: انتشار ۱۷، تحت کنترل
تغییر، پنجمین مجموعهٔ تغییرات فنی پس از نخستین نسخهٔ انتشار ۱۷، بدون به‌روزرسانی صرفاً ویرایشی پس از آن. \lr{TS 23.501 V2.0.0}:
پیش‌نویسی که برای تصویب به \acro{TSG} ارائه شده است (پیش از آنکه در انتشار ۱۵ به \lr{V15.0.0} تبدیل شود).
\end{examplebox}

\begin{examplebox}{بودجهٔ هم‌کنش‌پذیری}\label{ex:interop}
\textbf{مسئله.} یک تماس \acro{VoLTE} حدود ۲۰ واسط میان \acro{UE}، \acro{eNB}، \acro{MME}، دروازه‌ها، \acro{PCRF} و گره‌های
\acro{IMS} از فروشندگان مختلف را درگیر می‌کند. برای موفقیت سرتاسری ۹۹٪ چه درستی‌ای برای هر واسط لازم است؟

\textbf{حل.} $p \ge 0.99^{1/20} = 0.9995$. تجربهٔ میدانی \acro{VoLTE} اولیه، که به آزمون‌های گستردهٔ هم‌کنش‌پذیری و یک نمایهٔ
صنعتی برای محدودکردن گزینه‌ها نیاز داشت، این نکته را نشان می‌دهد.
\end{examplebox}

\begin{examplebox}{دسته‌بندی \lr{CR}ها}\label{ex:cr}
\textbf{مسئله.} دسته‌بندی کنید: (الف) مقدار نادرست یک زمان‌سنج در انتشار ۱۶ که در انتشار ۱۶ اصلاح می‌شود؛ (ب) همان اصلاح که
در انتشار ۱۷ کپی می‌شود؛ (ج) یک عنصر اطلاعاتی اختیاری جدید که یک قابلیت را ممکن می‌کند؛ (د) یک غلط تایپی.

\textbf{حل.} (الف) \lr{F}؛ (ب) \lr{A}؛ (ج) \lr{B} (فقط پیش از انجماد انتشار)؛ (د) \lr{D}.
\end{examplebox}

% ---------------------------------------------------------------------
\section{مقایسهٔ فناوری‌ها}
\begin{table}[htbp]
\centering\small
\caption{نهادهای استاندارد مرتبط با شبکه‌های موبایل.}\label{tab:bodies}
\rowcolors{2}{acnmist}{white}
\begin{tabularx}{\linewidth}{@{}l Y Y@{}}
\toprule
\rowcolor{white}نهاد & نقش & خروجی نمونه \\
\midrule
\lr{ITU-R}  & تنظیم طیف و نیازمندی‌های \acro{IMT}            & توصیه‌نامه‌ها و گزارش‌ها (مثلاً \lr{M.2083, M.2410, M.2160}) \\
\acro{3GPP} & مشخص‌کردن سامانه‌های موبایل                     & \acro{TS} و \acro{TR} (سری‌های ۲۱ تا ۳۸) \\
\acro{IETF} & پروتکل‌های اینترنت                              & \acro{RFC}ها (مثلاً \acro{SCTP}، \lr{Diameter}، \lr{HTTP/2}) \\
\acro{ETSI} & استانداردهای اروپایی؛ گروه‌های صنعتی \acro{NFV} و \acro{MEC} & \lr{EN}ها، مشخصات گروهی (مثلاً \acro{NFV}) \\
\lr{ONF, O-RAN Alliance} & واسط‌های باز برای \acro{SDN} و \acro{RAN} & مشخصات (مثلاً \lr{OpenFlow}، واسط‌های \lr{O-RAN}) \\
\acro{GSMA} & انجمن اپراتورها                                 & اسناد مرجع دائمی، نمایه‌ها \\
\bottomrule
\end{tabularx}
\end{table}

\paragraph{اجماع در برابر سرعت} استانداردسازی باز و مبتنی بر اجماع مشخصاتی تولید می‌کند که شرکت‌های بسیاری می‌توانند پیاده
کنند و برای دهه‌ها دوام دارند، به بهای زمان و پیچیدگی: گزینه‌ها برای راضی‌کردن طرف‌های بسیار افزوده می‌شوند. اتحادیه‌های
صنعتی و پروژه‌های متن‌باز سریع‌تر حرکت می‌کنند اما با تعهد رسمی کمتر؛ شبکه‌های امروزی همهٔ آن‌ها را ترکیب می‌کنند.

% ---------------------------------------------------------------------
\section{نکته‌های کلیدی}
\begin{takeawaybox}
\begin{itemize}
  \item \acro{3GPP} = هفت شریک سازمانی؛ کار در \lr{TSG RAN}، \lr{SA} و \lr{CT} و گروه‌های کاری آن‌ها؛ تصمیم‌های اجماعی.
  \item مرحلهٔ ۱ (نیازمندی‌ها)، مرحلهٔ ۲ (معماری)، مرحلهٔ ۳ (پروتکل‌ها)؛ \acro{TS} الزام‌آور، \acro{TR} آگاهی‌بخش.
  \item شمارهٔ سری موضوع را مشخص می‌کند (۲۳ معماری، ۲۴/۲۹ پروتکل‌ها، ۳۳ امنیت، ۳۶ \acro{LTE}، ۳۸ \acro{NR}).
  \item انتشارها مرحله‌به‌مرحله منجمد می‌شوند؛ نسخهٔ \lr{x.y.z} = انتشار، تغییرات فنی، تغییرات ویرایشی.
  \item درخواست‌های تغییر (\lr{F, A, B, C, D}) تنها راه تغییر مشخصه‌ای تحت کنترل تغییرند.
  \item هم‌کنش‌پذیری به‌صورت $p^k$ کاهش می‌یابد؛ مشخصات و آزمون باید دقیق باشند.
\end{itemize}
\end{takeawaybox}

% ---------------------------------------------------------------------
\section{تمرین‌ها}

\begin{exercisebox}{1}{کدام سری؟}\label{exr:series}
در کدام سری به دنبال این‌ها می‌گردید: (الف) پروتکل \acro{NAS} نسل پنجم؛ (ب) معماری امنیتی نسل پنجم؛ (ج) پروتکل \acro{MAC}
در \acro{LTE}؛ (د) نیازمندی‌های خدمت نسل پنجم؛ (هـ) پروتکل \lr{GTP-U}؟
\end{exercisebox}

\begin{exercisebox}{1}{\lr{TS} یا \lr{TR}؟}\label{exr:tstr}
آیا مدل کانالی که برای ارزیابی فناوری‌های نامزد به کار می‌رود الزام‌آور است؟ کجا منتشر می‌شود و چرا برای مقایسهٔ نتایج
شبیه‌سازی اهمیت دارد؟
\end{exercisebox}

\begin{exercisebox}{2}{نسخه‌ها}\label{exr:version}
مستندات یک فروشنده انطباق با «\lr{TS 24.501 V16.5.0}» را اعلام می‌کند. شبکهٔ اپراتور \lr{V16.9.0} را پیاده کرده است. چه
نوع تفاوت‌هایی ممکن است میان این دو وجود داشته باشد و چگونه آن‌ها را می‌یابید؟
\end{exercisebox}

\begin{exercisebox}{2}{واسط‌ها و درستی}\label{exr:interop}
با \cref{eq:interop}، $P_{\mathrm{e2e}}$ را برای $p = 0.995$ و $k = 20$ و ۵۰ محاسبه کنید. اگر $P_{\mathrm{e2e}} \ge 0.95$ لازم باشد،
یک خدمت با $p = 0.999$ از چند واسط می‌تواند بگذرد؟
\end{exercisebox}

\begin{exercisebox}{3}{انباشت}\label{exr:backlog}
یک گروه کاری ماهانه ۲۵۰ \acro{CR} را با زمان میانگین تصویب ۴ ماه پردازش می‌کند. یک انتشار جدید انفجاری می‌افزاید که ورودها
را به ۳۵۰ در ماه می‌رساند در حالی که ظرفیت تصویب ۲۵۰ در ماه می‌ماند. با \cref{eq:little} و یک استدلال سیالی ساده، انباشت
را پس از ۶ ماه انفجار تخمین بزنید. گروه چه می‌تواند بکند؟
\end{exercisebox}

\begin{exercisebox}{3}{مراحل هم‌پوشان}\label{exr:ttm}
با طول مراحل \cref{sec:ttm}، اگر پیاده‌سازی با تکمیل مرحلهٔ ۲ آغاز شود (هم‌پوشانی کامل با مرحلهٔ ۳)، محصولات تجاری چقدر زودتر
ظاهر می‌شوند؟ این چه خطری ایجاد می‌کند و تاریخ‌های انجماد چگونه آن را کاهش می‌دهند؟
\end{exercisebox}

\subsection*{پاسخ‌ها}
\addcontentsline{toc}{subsection}{پاسخ‌ها}

\begin{solutionbox}{\ref{exr:series}}
(الف) ۲۴ (\lr{TS 24.501})؛ (ب) ۳۳ (\lr{TS 33.501})؛ (ج) ۳۶ (\lr{TS 36.321})؛ (د) ۲۲ (\lr{TS 22.261})؛ (هـ) ۲۹ (\lr{TS 29.281}).
\end{solutionbox}

\begin{solutionbox}{\ref{exr:tstr}}
خیر: مدل‌های کانال آگاهی‌بخش‌اند و در \acro{TR}ها منتشر می‌شوند (مثلاً \lr{TR 38.901}). اهمیت دارند چون نتایج شبیه‌سازی
شرکت‌های مختلف فقط وقتی قابل مقایسه‌اند که از همان مدل کانال و فرض‌ها استفاده کنند، و \acro{TR} آن‌ها را تثبیت می‌کند.
\end{solutionbox}

\begin{solutionbox}{\ref{exr:version}}
هر دو انتشار ۱۶اند؛ \lr{V16.9.0} چهار مجموعهٔ تغییرات فنی بیشتر دارد -- بیشتر اصلاحات (دستهٔ \lr{F})، و احتمالاً شفاف‌سازی‌هایی
که رفتار را در موارد مرزی تغییر می‌دهند. فهرست \acro{CR}های میان دو نسخه را (پیوست تاریخچهٔ تغییرات و پایگاه دادهٔ \acro{CR})
مقایسه و بررسی کنید آیا هیچ‌یک بر رویه‌های پیاده‌شده اثر دارد.
\end{solutionbox}

\begin{solutionbox}{\ref{exr:interop}}
$0.995^{20} = 0.905$؛ $0.995^{50} = 0.778$. با $p = 0.999$: $k \le \ln 0.95/\ln 0.999 = 51.3$، یعنی حداکثر ۵۱ واسط.
\end{solutionbox}

\begin{solutionbox}{\ref{exr:backlog}}
در آغاز $N = 250\times4 = 1000$ \acro{CR} باز است. در طول انفجار انباشت ماهانه $350 - 250 = 100$ افزایش می‌یابد: پس از ۶ ماه
حدود ۱۶۰۰، و بنا به قانون لیتل زمان تا تصویب به حدود $1600/250 = 6.4$ ماه می‌رسد. گزینه‌ها: زمان جلسهٔ بیشتر (نشست‌های
الکترونیکی اضافی)، اولویت‌دادن به \acro{CR}های ضروری، به‌تعویق‌انداختن قابلیت‌های اختیاری به انتشار بعدی.
\end{solutionbox}

\begin{solutionbox}{\ref{exr:ttm}}
پیاده‌سازی به‌جای ۶ ماه با کل ۱۲ ماه مرحلهٔ ۳ هم‌پوشانی دارد: محصولات حدود ۶ ماه زودتر (۴۸ به‌جای ۵۴ ماه). خطر: پیاده‌سازی
یک هدف متحرک -- تغییرات بعدی مرحلهٔ ۳ دوباره‌کاری تحمیل می‌کنند و محصولات اولیه ممکن است هم‌کنش‌پذیر نباشند. تاریخ‌های
انجماد (به‌ویژه انجماد \lr{ASN.1}) نشان می‌دهند مشخصه کی برای پیاده‌سازی بدون چنین خطری به اندازهٔ کافی پایدار است.
\end{solutionbox}

% ---------------------------------------------------------------------
\section{مطالعهٔ بیشتر}
\lr{TR 21.900} \cite{tr21900} روش‌های کار \acro{3GPP}، شماره‌گذاری نسخه و دسته‌های \acro{CR} را توصیف می‌کند؛ \lr{TR 21.905}
\cite{tr21905} واژگان است. اسناد \lr{ITU-R} دربارهٔ \lr{IMT-2020} و \lr{IMT-2030} \cite{itu2015m2083,itu2017m2410,itu2023m2160}
نشان می‌دهند کار \acro{3GPP} چگونه به نیازمندی‌های جهانی پیوند می‌خورد. دالمن و همکاران \cite{dahlman2020nr} تاریخچهٔ
استانداردسازی \acro{NR} را توصیف می‌کنند.

\begin{latin}
\printbibliography[heading=subbibliography,title={References}]
\end{latin}
\end{document}
